\documentclass[french]{article}

\usepackage[utf8]{inputenc}
\usepackage[T1]{fontenc}
\usepackage{lmodern}
\usepackage{geometry}
\usepackage{graphicx}
\usepackage{babel}
\usepackage{amsmath}
\usepackage{amsthm}
\usepackage{amssymb}
\usepackage{hyperref}
\usepackage{stmaryrd}
\usepackage{enumitem}
\usepackage[most]{tcolorbox}
\usepackage{dsfont}
\usepackage{multirow}
\usepackage{etoc}
\usepackage{titlesec}
\usepackage{esint}
\usepackage{cancel}
\usepackage{mathdots}

\setlist[itemize]{label=$\bullet$}


\renewcommand{\P}{\mathbb{P}}
\newcommand{\N}{\mathbb{N}}
\newcommand{\Z}{\mathbb{Z}}
\newcommand{\Q}{\mathbb{Q}}
\newcommand{\R}{\mathbb{R}}
\newcommand{\C}{\mathbb{C}}
\newcommand{\K}{\mathbb{K}}
\renewcommand{\L}{\mathbb{L}}
\newcommand{\U}{\mathbb{U}}
\newcommand{\st}{^*}
\newcommand{\tq}{\middle|}
\newcommand{\inv}{^{-1}}
\newcommand{\E}{\mathbb{E}}
\newcommand{\F}{\mathbb{F}}
\newcommand{\V}{\mathbb{V}}
\renewcommand{\ker}{\text{Ker}}
\newcommand{\im}{\text{Im}}
\newcommand{\card}{\text{Card}}
\newcommand{\Tr}{\text{Tr}}
\newcommand{\Vect}{\text{Vect}}
\newcommand{\rg}{\text{rg}}

\theoremstyle{plain}
\newtheorem*{ind}{Indication}

\theoremstyle{definition}

\newtheorem*{ex}{Exemple}
\newtheorem*{met}{Point méthode}
\newtheorem{exo}{Exercice}
\newtheorem*{rmq}{Remarque}

\newtcolorbox[auto counter]{dfn}{
	enhanced,
	breakable,
	colback=white,
	fonttitle=\sc,
	title={Définition \thetcbcounter}
}

\newtcolorbox[auto counter]{thm}[2][]{enhanced, breakable, colback=white, fonttitle=\sc, title=Théorème~\thetcbcounter~: #2}

\theoremstyle{remark}
\newtheorem*{app}{Application}

\title{Topologie des espaces vectoriels normés}
\date{}

\begin{document}
\tableofcontents
%"Le  programme est plutôt succinct, il manque des choses"

\newpage
\maketitle

Dans tout le chapitre, la lettre $\K$ désigne le corps $\R$ ou le corps $\C$. Tous les espaces vectoriels considérés sont des espaces vectoriels réels ou complexes.

\section{Généralités}

\subsection{Définition}

\begin{dfn}
	Étant donné un $\K$-espace vectoriel $E$, on appelle \textbf{norme sur $E$} toute application $N$ de $E$ dans $\R_+$ vérifiant les trois propriétés suivantes :
	\begin{itemize}
		\item \textbf{homogénéité} : $\forall (\lambda,x)\in\K\times E,\ N(\lambda x)=\lvert\lambda\rvert N(x)$ ;
		\item \textbf{inégalité triangulaire} : $\forall (x,y)\in E^2,\ N(x+y)\leqslant N(x)+N(y)$ ;
		\item \textbf{séparation} : $\forall x\in E,\ N(x)=0\implies x=0$.
	\end{itemize}
	Tout espace vectoriel muni d'une norme est appelé \textbf{espace vectoriel normé}.
\end{dfn}

\paragraph{Notations}
\begin{itemize}
	\item On note $(E,N)$ pour désigner l'espace vectoriel $E$ muni d'une norme $N$. Lorsqu'il n'y a pas de risque d'ambiguïté quant à la norme utilisée, on peut ne pas la préciser et se contenter de noter $E$.
	\item La plupart du temps, il est d'usage de noter $\lVert \cdot\rVert$ la norme utilisée, et donc $\lVert x\rVert$ la norme d'un vecteur $x$.
\end{itemize}

\begin{rmq}
	\begin{itemize}
		\item Si $N$ est une norme, alors la propriété d'homogénéité implique que $N(0)=0$. Dans la définition précédente, l'implication utilisée pour formuler la propriété de séparation est donc une équivalence.
		\item Une application de $E$ dans $\R_+$ qui vérifie l'inégalité triangulaire et l'homogénéité est appelée une \textbf{semi-norme} (HP).
	\end{itemize}
\end{rmq}

\begin{ex}
	\begin{enumerate}
		\item L'application $x\mapsto\lvert x\rvert$ est une norme sur $\K$ (où $\lvert\cdot\rvert$ désigne la valeur absolue sir $\K=\R$ et le module sir $\K=\C$).
		\item Les applications : $$\begin{array}{rcl}
			\K&\longrightarrow&\R_+\\
			(x,y)&\longmapsto&\max(\lvert x\rvert,\lvert y\rvert)
		\end{array}\ \ \text{et}\ \ 
		\begin{array}{rcl}
			\K&\longrightarrow&\R_+\\
			(x,y)&\longmapsto&\lvert x\rvert+\lvert y\rvert
		\end{array}$$ sont des normes sur $\K^2$, appelées respectivement \textbf{norme infinie} et \textbf{norme un}, et notées $\lVert\cdot\rVert_\infty$ et $\lVert\cdot\rVert_1$.
	\end{enumerate}
\end{ex}

\begin{thm}{}%prop
	Soit $E$ un espace préhilbertien réel, c'est-à-dire un $\R$-espace vectoriel muni d'un produit scalaire $(\cdot\mid\cdot)$. L'application $x\mapsto\sqrt{(x\mid x)}$ est une norme sur $E$, appelée \textbf{norme euclidienne} associée au produit scalaire $(\cdot\mid\cdot)$.
\end{thm}

\begin{ex}
	De la proposition précédente, on déduit que les espaces préhilbertiens réels sont naturellement munis d'une structure d'espace vectoriel normé :
	\begin{enumerate}
		\item l'espace $\R^n$ muni du produit scalaire canonique $$\left((x_1,\ldots,x_n)\mid(y_1,\ldots,y_n)\right)=\sum_{k=1}^{n}x_ky_k$$ dont la norme associée est donnée par $\lVert(x_1,\ldots,x_n)\rVert=\sqrt{x_1^2+\cdots+x_n^2}$ ;
		\item l'espace $\mathcal{C}^0([a,b],\R)$ (avec $a<b$) muni du produit scalaire $(f\mid g)=\displaystyle\int_{a}^{b}fg$ dont la norme associée est donnée par $\lVert f\rVert=\displaystyle\sqrt{\int_{a}^{b}f^2}$.
	\end{enumerate}
\end{ex}

\paragraph{Terminologie} Dans chacun des deux exemples précédents, la norme euclidienne considérée est appelée \textbf{norme deux}. Tout cela sera précisé dans la section "Exemples d'espaces vectoriels normés".

\subsubsection*{Vecteur unitaire}

\begin{dfn}
	On dit qu'un vecteur $x$ d'un espace vectoriel normé $E$ est \textbf{unitaire} lorsque $\lVert x\rVert=1$.
\end{dfn}

\begin{exo}
	On suppose que $\K=\R$. Soit $x$ un vecteur non nul d'un espace vectoriel normé $E$. Montrer qu'il existe un unique vecteur unitaire colinéaire à $x$ et de même sens que lui.
\end{exo}

\paragraph{Terminologie} Étant donné un vecteur $x$ non nul de $E$, le vecteur $\displaystyle\frac{x}{\lVert x\rVert}$ est appelé \textbf{vecteur unitaire associé à $x$}.

\begin{exo}
	Donner une condition nécessaire et suffisante sur $E$ pour qu'il existe dans $E$ des vecteurs unitaires.
\end{exo}


\subsection{Inégalités triangulaires}

Nous avons défini l'inégalité triangulaire (propriété vérifiée par toute norme par définition) de la manière suivante : $$\forall (x,y)\in E^2,\ \lVert x+y\rVert\leqslant\lVert x\rVert+\lVert y\rVert$$ De manière plus complète, nous regroupons sous le nom d'\textit{inégalités triangulaires} les inégalités du résultat suivant :

\begin{thm}{Inégalités triangulaires}%prop
	Soit $x$ et $y$ deux vecteurs d'un espace vectoriel normée $E$. Alors on a : $$\lvert \lVert x\rVert-\lVert y\rVert\rvert\leqslant\lVert x\pm y\rVert\leqslant\lVert x\rVert+\lVert y\rVert$$
\end{thm}

\paragraph{Terminologie} L'inégalité $\lvert \lVert x\rVert-\lVert y\rVert\rvert\leqslant\lVert x\pm y\rVert$ est parfois appelée \textit{seconde inégalité triangulaire}.

\begin{exo}
	\begin{enumerate}
		\item Montrer que si $x$ et $y$ sont positivement colinéaires, alors $\lVert x+y\rVert=\lVert x\rVert+\lVert y\rVert$.
		\item On se place dans $\R^2$ muni de la norme un. Exhiber deux vecteurs $u_1$ et $u_2$ non colinéaires tels que $\lVert u_1+u_2\rVert=\lVert u_1\rVert+\lVert u_2\rVert$.
	\end{enumerate}
\end{exo}

\paragraph{Attention}\begin{itemize}
	\item Contrairement à l'intuition géométrique, et comme le montre l'exercice précédent, dans le cas d'égalité de l'inégalité triangulaire, on ne peut, \textit{a priori}, rien en déduire.
	\item En revanche, si la norme considérée est euclidienne, c'est-à-dire si elle est associée à un produit scalaire, alors le cas d'égalité dans l'inégalité triangulaire est celui où les deux vecteurs sont positivement colinéaires.
\end{itemize}

\begin{exo}
	Démontrer l'affirmation précédente.
\end{exo}

\paragraph{Inégalité triangulaire généralisée} La propriété d'homogénéité combinée à l'inégalité triangulaire permet, par récurrence, de montrer que si $x_1,\ldots,x_n$ sont des éléments de $E$ et $\lambda_1,\ldots,\lambda_n$ sont des scalaires, alors on a : $$\left\lVert\sum_{k=1}^{n}\lambda_kx_k\right\rVert\leqslant\sum_{k=1}^{n}\lvert\lambda_k\rvert\lVert x_k\rVert$$ Dorénavant, nous utiliserons librement ce résultat.

\begin{exo}
	Montrer réciproquement qu'une application $N:E\to\R_+$ qui vérifie : $$\forall (x,y)\in E^2,\ \forall (\lambda,\mu)\in\K^2,\ \lVert\lambda x+\mu y\rVert\leqslant\lvert\lambda\rvert\lVert x\rVert+\lvert\mu\rvert\lVert y\rVert$$ est une semi-norme.
\end{exo}

\subsection{Distance associée à une norme}

\begin{dfn}
	On appelle \textbf{distance associée à la norme} $\lVert\cdot\rVert$ l'application : $$\begin{array}{lrcl}
		d:&E^2&\longrightarrow&\R\\
		&(x,y)&\longmapsto&\lVert x-y\rVert
	\end{array}$$
\end{dfn}

\begin{exo}
	Montrer que la distance $d$ associée à la norme $\lVert\cdot\rVert$ vérifié les propriétés suivantes :
	\begin{itemize}
		\item \textbf{séparation} : $\forall (M,N)\in E^2,\ d(M,N)=0\iff M=N$ ;
		\item \textbf{symétrie} : $\forall(M,N)\in E^2,\ d(M,N)=d(N,M)$ ;
		\item \textbf{inégalité triangulaire} : $\forall (M,N,P)\in E^3,\ d(M,P)\leqslant d(M,N)+d(N,P)$ ;
		\item \textbf{invariance par translation} : $\forall (M,N,u)\in E^3,\ d(M+u,N+u)=d(M,N)$.
	\end{itemize}
\end{exo}

\subsubsection*{Formulation des inégalités triangulaires en termes de distance}

\begin{thm}{}%prop
	Si $M$, $N$ et $P$ désignent trois points de $E$, alors on a : $$\lvert d(M,N)-d(N,P)\rvert\leqslant d(M,P)\leqslant d(M,N)+d(N,P)$$
\end{thm}

De manière moins formelle, l'inégalité $d(M,P)\leqslant d(M,N)+d(N,P)$ signifie que dans un espace vectoriel normé, les propriétés "géométriquement intuitives" suivantes sont vraies :
\begin{itemize}
	\item dans un triangle, la longueur d'un côté est inférieure à la somme des longueurs des deux autres ;
	\item \textit{un} plus court chemin entre deux points est la ligne droite.
\end{itemize}

\paragraph{Attention} Bien comprendre la raison d'un choix de l'article indéfini \textit{un} dans la phrase "\textit{un} plus court chemin entre deux points" : la ligne droite est un plus court chemin entre deux points, mais ce n'est pas nécessairement le seul (du moins si la norme n'est pas euclidienne ; produit un exemple).

\subsection{Espaces métriques (HP)}

\begin{dfn}
	Soit $A$ un ensemble. Une application $d:A^2\to\R_+$ est une \textbf{distance} sur $A$ si elle vérifie les trois axiomes :
	\begin{enumerate}[label=(\arabic*)]
		\item \textbf{séparation} : $\forall (x,y)\in A^2,\ d(x,y)=0\iff x=y$ ;
		\item \textbf{symétrie} : $\forall(,y)\in A^2,\ d(x,y)=d(y,x)$ ;
		\item \textbf{inégalité triangulaire} : $\forall (x,y,z)\in A^3,\ d(x,z)\leqslant d(x,y)+d(y,z)$.
	\end{enumerate}
	Un ensemble $A$ muni d'une distance est appelé un \textbf{espace métrique}.
\end{dfn}

\begin{rmq}
	\begin{itemize}
		\item Si on remplace $(1)$ par $d(x,x)=0$ (suppression de la séparation), $d$ est alors appelé une \textbf{semi-distance} sur $A$.
		\item Des axiomes $(2)$ et $(3)$ on tire l'inégalité : $$d(x,z)\geqslant\lvert d(x,y)-d(y,z)\rvert$$ d'où la propriété : "dans un triangle, la longueur de tout côté est comprise entre la somme et la différence des longueurs des deux autres".
	\end{itemize}
\end{rmq}

\begin{ex}
	Sur $\K$, la distance usuelle $(x,y)\mapsto\lvert x-y\rvert$ associée à la norme $\lvert\cdot\rvert$.
\end{ex}

Plus généralement :

\begin{ex}%ex fond
	Si $(E,N)$ est un espace vectoriel normé, et si $A$ est une partie de $E$, l'application $(x,y)\mapsto\lVert x-y\rVert$ définit une distance sur $A$. \\
	En particulier, $E$ muni de la distance associée à $N$ est un espace métrique.
\end{ex}

Mais il existe des espaces métriques plus exotiques dont la distance n'est pas directement associée à une norme.

\begin{ex}
	Sur n'importe quel ensemble $A$, la \textbf{distance discrète} définie par : $$d(x,y)=\begin{cases}
		0&\text{si }x=y\\
		1&\text{si }x\ne y 
	\end{cases}$$ fait de $A$ un espace métrique.
\end{ex}

\begin{ex}
	Soit $d$ une distance sur un ensemble $A$. Montrer que $\delta:(x,y)\mapsto\min(1,d(x,y))$ définit également une distance sur $A$.
\end{ex}

\begin{ex}
	Dans $\C$, on définit $d(a,b)=\lvert b-a\rvert$ sir $a$ et $b$ sont positivement colinéaires, et $d(a,b)=\lvert a\rvert+\lvert b\rvert$ sinon. $d$ est une distance sur $\C$, appelée \textbf{distance SNCF}.
\end{ex}

\begin{rmq}%rmq importante
	Toutes les définitions et propriétés qui vont suivre seront formulées dans le cas d'un espace métrique non vide $(E,d)$. Pour rester dans le programme, il suffit de considérer un espace vectoriel normé $E$ et sa distance $d$ associée.
\end{rmq}

\subsection{Boule fermée, boule ouverte}

\begin{dfn}
	Soit $a$ un élément de $E$ et $r$ un réel strictement positif.
	\begin{itemize}
		\item On appelle \textbf{boule ouverte de centre $a$ et de rayon $r$}, et l'on note $B_o(a,r)$ la partie : $$B_o(a,r)=\{x\in E\ : \ d(a,x)<r\}$$
		\item On appelle \textbf{boule fermée de centre $a$ et de rayon $r$}, et l'on note $B_f(a,r)$ la partie : $$B_f(a,r)=\{x\in E\ : \ d(a,x)\leqslant r\}$$
		\item On appelle \textbf{sphère de centre $a$ et de rayon $r$}, et l'on note $S(a,r)$ la partie : $$S(a,r)=\{x\in E\ : \ d(a,x)=r\}=B_f(a,r)\setminus B_o(a,r)$$
	\end{itemize}
\end{dfn}

\paragraph{Notation} Dans $\C$, on parle plutôt de \textit{disque}, et l'on utilise alors les notations $D_o(a,r)$ et $D_f(a,r)$ au lieu de $B_o(a,r)$ et $B_f(a,r)$ respectivement.

\begin{rmq}
	Il est clair, d'après la définition, que les boules dépendant de la distance considérée sur $E$ : une partie de $E$ peut être une boule pour une distance et ne pas l'être pour une autre distance.
\end{rmq}

\begin{ex}
	Boules et sphères pour la distance discrète.
\end{ex}

\begin{exo}
	Quelles sont les boules fermés pour la distance SNCF ? Expliquer cette appellation.
\end{exo}

\subsection{Convexité dans un espace vectoriel normé}

Dans cette section, $E$ est un espace vectoriel réel ou complexe.

\subsubsection*{Parties convexes}

\begin{dfn}
	Une partie $A$ de $E$ est dite \textbf{convexe} si : $$\forall (x,y)\in A^2,\ \forall \lambda\in[0,1],\ (1-\lambda)x+\lambda y\in A$$
\end{dfn}

\begin{ex}
	\begin{enumerate}
		\item L'ensemble vide est convexe.
		\item Tout singleton est convexe.
		\item Tout sous-espace affine de $E$ est convexe.
		\item Les parties convexes de $\R$ sont les intervalles.
	\end{enumerate}
\end{ex}

\paragraph{Terminologie} Étant donnés deux points $x$ et $y$,  l'ensemble des points de la forme $$(1-\lambda)x+\lambda y\ \ \text{avec}\ \ \lambda\in[0,1]$$ est appelé \textbf{segment} $[x,y]$.

\paragraph{Interprétation géométrique} Dire qu'une partie $A$ est convexe signifie qu'étant donnés deux points de $A$, le segment reliant ces deux points est inclus dans $A$.

\begin{thm}{Stabilité par passage au barycentre à coefficients positifs (HP)}%prop
	Si $A$ est une partie convexe d'un espace vectoriel réel $E$, alors on a : $$\forall p\in\N\st,\ \forall (x_1,\ldots,x_p)\in A^p,\ \forall(\lambda_1,\ldots,\lambda_p)\in\R_+^p,\ \sum_{k=1}^{p}\lambda_k=1\implies\sum_{k=1}^{p}\lambda_kx_k\in A$$
\end{thm}
\begin{proof}
	Procéder par récurrence.
\end{proof}

\begin{ex}
	Soit $C$ un convexe de $\C$. Si $X$ est une variable aléatoire finie à valeurs dans $C$, son espérance est également dans $C$.
\end{ex}

\subsubsection*{Boules d'un espace vectoriel normé}

\begin{thm}{}%prop
	Toute boule (ouvert ou fermée) d'un espace vectoriel normé est une partie convexe.
\end{thm}

\begin{rmq}
	Dans un espace vectoriel normé, la connaissance de la boule unité fermée $B_f(0,1)$ suffit pour connaître toutes les boules fermées. En effet, si $a$ est un élément de $E$ et $r$ un réel strictement positif, la boule $B_f(a,r)$ est l'image de $B_f(0,1)$ par... \\
	Toutes les boules ont donc "la même forme".
\end{rmq}

\begin{ex}
	Sur un même plan, dessin des boules unité de $\R^2$ pour les normes infinie, un et deux.
\end{ex}

Exercice Moulin.EVN.13

\subsection{Exemples d'espaces vectoriels normés}

Le résultat suivant sera utile pour prouver l'homogénéité d'une normé définie par une borne supérieure.

%"Dans sa grande sagesse, le programme nous donne ce lemme. Pour une fois que je peux mettre sagesse et prgramme dans la même phrase...

\begin{thm}{}%lem
	Étant donnée $A$ une partie non vide de $\R$ et $k\in\R_+$, on a : $$\sup(kA)=k\sup(A)$$
\end{thm}

\begin{rmq}
	Ce résultat s'applique en particulier si la partie $A$ admet un plus grand élément, et l'on a dans ce cas $\max(kA)=k\max(A)$.
\end{rmq}

\subsubsection*{L'espace $\K^n$}

Soit $x$ un vecteur de $\K^n$ dont on note $(x_1,\ldots,x_n)$ les composantes. \\
L'espace $\K^n$ peut être peut être muni d'une structure d'espace vectoriel normé à l'aide de différentes normes, dont les trois normes classiques suivantes :
\begin{itemize}
	\item la \textbf{norme infinie}, définie par : $$\lVert x\rVert_\infty=\max(\lvert x_1\rvert,\ldots,\lvert x_n\rvert)$$
	\item la \textbf{norme un}, définie par : $$\lVert x\rVert_1=\lvert x_1\rvert+\cdots+\lvert x_n\rvert$$
	\item la \textbf{norme deux}, ou \textbf{norme euclidienne canonique}, définie par : $$\lVert x\rVert_2=\sqrt{\lvert x_1\rvert^2+\cdots+\lvert x_n\rvert^2}$$
\end{itemize}

\begin{thm}{}
	Les trois applications :
	$$\begin{array}{rcl}
		\K^n&\longrightarrow&\R_+\\
		x&\longmapsto&\lVert x\rVert_\infty
	\end{array}\ \
	\begin{array}{rcl}
		\K^n&\longrightarrow&\R_+\\
		x&\longmapsto&\lVert x\rVert_1
	\end{array}\ \
	\begin{array}{rcl}
		\K^n&\longrightarrow&\R_+\\
		x&\longmapsto&\lVert x\rVert_2
	\end{array}$$ sont des normes. Chacune d'entre elles munit donc $\K^n$ d'une structure d'espace vectoriel normé.
\end{thm}

\begin{rmq}
	Si $\K=\R$, alors on peut aussi écrire $\lVert x\rVert_2=\sqrt{x_1^2+\cdots+x_n^2}$.
\end{rmq}

\paragraph{Approfondissement : norme $p$} Pour $p\geqslant 1$, il arrive parfois de munir $\K^n$ de la norme suivante, appelée \textbf{norme $p$}, définie ainsi : $$\lVert x\rVert_p=(\lvert x_1\rvert^p+\cdots+\lvert x_n\rvert^p)^{1/p}$$ Pour plus de détails, voir l'exercice Moulin.EVN.6.

\subsubsection*{Généralisation à tout espace vectoriel muni d'une base}
Plus généralement, si $E$ est un $\K$-espace vectoriel normé non nul muni d'une base $\mathcal{B}=(u_i)_{i\in I}$, alors on dispose sur $E$ des trois normes classiques suivantes, appelées respectivement, \textbf{norme infinie}, \textbf{norme un} et \textbf{norme deux} : $$\lVert x\rVert=\underset{i\in I}{\max}\lvert x_i\rvert\ ; \ \lVert x\rVert_1=\sum_{i\in I}\lvert x_i\rvert\ \text{et}\ \lVert x\rVert_2=\sqrt{\sum_{i\in I}\lvert x_i\rvert^2}$$ où, étant donné un élément $x$ de $E$, la famille $(x_i)_{i\in I}$ désigne les composantes de $x$ dans la base $\mathcal{B}$.

\begin{rmq}
	Les définitions ci-dessus des normes un, deux et infinie sont licites car, par définition, la famille $(x_i)_{i\in I}$ des composantes de $x$ dans la base $\mathcal{B}$ est à support fini (c'est-à-dire ne comporte qu'un nombre fini de termes non nuls).
\end{rmq}

\subsubsection*{Espace des fonctions bornées}
L'ensemble $\mathcal{B}(X,\K)$ des fonctions bornées d'un ensemble $X$ non vide à valeurs dans $\K$ est un sous-espace vectoriel du $\K$-espace vectoriel $\mathcal{F}(X,\K)$ des applications de $X$ dans $\K$.

\begin{thm}{}
	L'application : $$\begin{array}{lrcl}
		\mathcal{N}_\infty&\mathcal{B}(X,\K)&\longrightarrow&\R_+\\
		&f&\longmapsto&\underset{X}{\sup}\lvert f\rvert
	\end{array}$$ est une norme sur $\mathcal{B}(X,\K)$. On l'appelle \textbf{norme infinie}.
\end{thm}
La norme infinie munit donc $\mathcal{B}(X,\K)$ d'une structure d'espace vectoriel normé.

\begin{ex}
	Cela s'applique en particulier à l'espace vectoriel $\mathcal{B}(\N,\K)$ des suites bornées à valeurs dans $\K$. Pour $(u_n)_{n\in\N}\in\K^\N$, on a : $$\mathcal{N}_\infty(u)=\underset{\N}{\sup}\lvert u_n\rvert$$
\end{ex}

\subsubsection*{Espace des fonctions continues sur un segment}

Soit $a$ et $b$ deux réels tels que $a<b$. On note $\mathcal{C}([a,b],\K)$ l'ensemble des fonctions continues du segment $[a,b]$ dans $\K$. On sait que $\mathcal{C}([a,b],\K)$ est un sous-espace vectoriel de $\mathcal{F}([a,b],\K)$, et même de $\mathcal{B}([a,b],\K)$.\\
Cet espace peut être muni d'une structure d'espace vectoriel normé, en particulier à l'aide des deux normes classiques suivantes :
\begin{itemize}
	\item la \textbf{norme un}, définie par : $\mathcal{N}_1(f)=\displaystyle\int_{a}^{b}\lvert f\rvert$ ;
	\item la \textbf{norme deux}, définie par : $\mathcal{N}_2(f)=\displaystyle\sqrt{\int_{a}^{b}\lvert f\rvert^2}$
\end{itemize}

\begin{exo}
	Montrer que ces applications sont des normes sur $\mathcal{C}([a,b],\K)$. Sont-elles des normes sur $\mathcal{CM}([a,b],\K)$ ?
\end{exo}

\paragraph{Approfondissement : norme $p$} Pour $p\geqslant 1$, on munit parfois l'espace $\mathcal{C}([a,b],\K)$ de la norme suivante, appelée \textbf{norme $p$}, définie ainsi : $$\mathcal{N}_p(f)=\left(\int_{a}^{b}\lvert f\rvert^p\right)^{1/p}$$
Pour plus de détails, voir l'exercice Moulin.EVN.07.

\subsection{Obtention de nouveaux espaces vectoriels normés}

\subsubsection*{Norme induite}

Si $F$ est un sous-espace vectoriel de $E$, alors la norme $\lVert \cdot\rVert$ dont on dispose sur $E$ induit naturellement une structure d'espace vectoriel normé sur $F$, par sa restriction à $F$ : $$\begin{array}{rcl}
	F&\longrightarrow&\R_+\\
	x&\longmapsto&\lVert x\rVert.
\end{array}$$ Cette restriction est appelé \textbf{norme induite} sur $F$ par la norme $\lVert\cdot\rVert$.

\begin{ex}
	\begin{enumerate}
		\item La valeur absolue sur $\R$ est induite par le module sur $\C$.
		\item Soit $a$ et $b$ deux réels tels que $a<b$. Il est clair que $\mathcal{C}([a,b],\K)$ est un sous-espace vectoriel de $\mathcal{B}([a,b],\K)$, dont on a vu que la norme infinie le munissait d'une structure d'espace vectoriel normé. La norme induit associée : $$\begin{array}{rcl}
			\mathcal{C}([a,b],\K)&\longrightarrow&\R_+\\
			f&\longmapsto&\mathcal{N}_\infty(f)
		\end{array}$$ constitue alors, en plus des normes un et deux, une autre norme sur $\mathcal{C}([a,b],\K)$.
	\end{enumerate}
\end{ex}

\subsubsection*{Application à $\K[X]$}

En considérant $\K[X]$ comme l'espace vectoriel des suites presque nulles, on peut le munir des trois normes usuelles définies à l'aide des coordonnées dans la base canonique par :
$$\lVert P\rVert_1=\sum_{n=0}^{+\infty}\lvert a_n\rvert\ \ \lVert P\rVert=\sqrt{\sum_{n=0}^{+\infty}\lvert a_n\rvert^2}\ \ \text{et}\ \ \lVert P\rVert_\infty=\underset{n\in\N}{\sup}\lvert a_n\rvert$$

où $P=\displaystyle\sum_{n=0}^{+\infty}a_nX^n$, la suite $(a_n)_{n\in\N}$ étant nulle à partir d'un certain rang. En effet, ces trois normes sont naturellement définies sur des espaces de suites plus grands.\\
Mais on peut aussi considérer un polynôme comme une fonction continue sur un segment $[a,b]$ avec $a<b$, ce qui permet également de munir $\K[X]$ des trois normes usuelles : 

$$\mathcal{N}_1(P)=\int_{a}^{b}\lvert P\rvert\ \ \mathcal{N}_2(P)=\sqrt{\int_{a}^{b}\lvert P\rvert^2}\ \ \text{et}\ \ \mathcal{N}_\infty(P)=\underset{[a,b]}{\max}\lvert P\rvert$$

\begin{exo}
	Pourquoi s'agit-il effectivement de normes ?
\end{exo}

\subsubsection*{Utilisation d'une application linéaire injective}

Soit $E$ et $F$ deux $\K$-espaces vectoriels. Si l'on dispose d'une norme $\lVert\cdot\rVert$ sur $E$, et si $u$ est une application linéaire injective de $F$ dans $E$, alors l'application : $$\begin{array}{lrcl}
	N:&F&\longrightarrow&\R_+\\
	&x&\longmapsto&\lVert u(x)\rVert
\end{array}$$ est une norme sur $F$.

\begin{ex}
	\begin{enumerate}
		\item On retrouve ainsi les résultats de l'exercice précédent.
		\item Si $F$ est un sous-espace vectoriel de $E$, alors la norme induite sur $F$ par celle de $E$ n'est rien d'autre que la norme obtenue en utilisant l'injection canonique de $F$ dans $E$ comme application linéaire injective.
		\item Soi $(a,b)\in\R^2$ avec $a<b$. Notons $E=\mathcal{C}([a,b],\K)$.\\
		Si $\omega:[a,b]\to\K$ est une fonction continue et à valeurs strictement positives, alors l'application : $$\begin{array}{lrcl}
			N:&E&\longrightarrow&\R_+\\
			&f&\longmapsto&\displaystyle\int_{a}^{b}\omega(x)\lvert f(x)\rvert dx
		\end{array}$$ est une norme sur $E$. En effet, c'est la norme que l'on obtient avec la norme $1$ et l'application linéaire injective : $$\begin{array}{rcl}
		E&\longrightarrow&E\\
		f&\longmapsto&(x\mapsto\omega(x)f(x))
		\end{array}$$
	\end{enumerate}
\end{ex}

\subsubsection*{Produit fini d'espaces vectoriels normés}

\begin{thm}{}%prop
	Soit $p\in\N\st$ ainsi que $p$ espaces vectoriels normés $(E_1,\lVert\cdot\rVert_1),\ldots,(E_p,\lVert\cdot\rVert_p)$. Alors, l'application : $$\begin{array}{rcl}
		E_1\times\cdots\times E_p&\longrightarrow&\R_+\\
		(x_1,\ldots,x_p)&\longmapsto&\max(\lVert x_1\rVert_1,\ldots,\lVert x_p\rVert_p)
	\end{array}$$ est une norme, et munit donc l'espace produit $E_1\times\cdots\times E_p$ d'une structure d'espace vectoriel normé. Cette norme est appelée \textbf{norme produit}.
\end{thm}

\begin{ex}
	La norme infinie sur $\K^n$ est la norme produit obtenue en considérant, sur $\K$, la norme $x\mapsto\lvert x\rvert$.
\end{ex}

\begin{rmq}
	En plus de la norme produit présentée dans la proposition précédente, d'autres normes peuvent être considérées sur un espace produit, comme le montre l'exercice suivant.
\end{rmq}

\begin{exo}
	Avec les notations de la proposition précédente, montrer que les applications $N_1$ et $N_2$ définies par : 
	$$\begin{array}{lrcl}
		N_1:&E_1\times\cdots\times E_p&\longrightarrow&\R_+\\
		&(x_1,\ldots,x_p)&\longmapsto&\displaystyle\sum_{k=1}^{p}\lVert x_k\rVert_k
	\end{array}\ \ \text{et}\ \
	\begin{array}{lrcl}
		N_2:&E_1\times\cdots\times E_p&\longrightarrow&\R_+\\
		&(x_1,\ldots,x_p)&\longmapsto&\displaystyle\sqrt{\sum_{k=1}^{p}\lVert x_k\rVert_k^2}
	\end{array}$$ sont des normes sur l'espace produit $E_1\times\cdots\times E_p$.
\end{exo}

\begin{rmq}
	Plus généralement, si $(E_1,d_1),\ldots,(E_p,d_p)$ sont des espaces métriques, on peut définie la \textbf{distance produit} sur $E_1\times\cdots\times E_p$ en posant : $$d((x_1,\ldots,x_p),(y_1,\ldots,y_p))=\max(d_1(x_1,y_1),\ldots,d_p(x_p,y_p))$$
\end{rmq}

\begin{exo}
	Dans l'espace produit $E\times F$ :
	\begin{enumerate}
		\item Quelle est la boule fermée (respectivement ouverte) centrée en $(a,b)\in E\times F$ et de rayon $r>0$ ?
		\item Quelle est la sphère centrée en $(a,b)\in E\times F$ et de rayon $r>0$ ?
	\end{enumerate}
\end{exo}

\subsection{Distance à une partie}

\begin{dfn}
	Étant donnée une partie $A$ non vide de $E$ ainsi qu'un élément $x$ de $E$, on appelle la \textbf{distance de $x$ à $A$}, et on note $d(x,A)$, la quantité : $$d(x,A)=\inf\left\{d(x,a)\ \mid \ a\in A\right\}$$
\end{dfn}

\begin{rmq}
	Avec les notations de la définition précédente :
	\begin{enumerate}
		\item Le caractère non vide de $A$ implique que $\{d(x,a)\ \mid\ a\in A\}$ est une partie non vide de $\R$ minorée par $0$, ce qui assure l'existence de sa borne inférieure et le fait que celle-ci soit positive ;
		\item Dans le cas d'un espace vectoriel normé, par définition de la distance $d$ associée à la norme, on a aussi : $$d(x,A)=\inf\{\lVert x-a\rVert\ \mid\ a\in A\}$$ 
	\end{enumerate}
\end{rmq}

\begin{thm}{}%prop
	Soit $A$ une partie non vide de $E$ ainsi que $x$ et $y$ deux éléments de $E$. On a : $$\lvert d(x,A)-d(y,A)\rvert\leqslant d(x,y)$$
\end{thm}
\begin{proof}
	Commencer par démontrer l'inégalité $d(x,A)-d(y,A)\leqslant d(x,y)$ puis utiliser la symétrie en $x$ et $y$.
\end{proof}

\begin{rmq}
	Le résultat précédente se reformule ainsi : l'application $$\begin{array}{rcl}
		(E,\lVert\cdot\rVert)&\longrightarrow&(\R,\lvert\cdot\rvert)\\
		x&\longmapsto&d(x,A)
	\end{array}$$ est $1$-lipschitzienne.
\end{rmq}

\paragraph{Attention}
\begin{itemize}
	\item Si $x$ appartient à $A$, alors il est clair que $d(x,A)=0$. En revanche, la réciproque est fausse. Par exemple, dans $\R$ muni de la valeur absolue, on a $d(0,\R_+\st)=0$ et pourtant $0\notin\R_+\st$.
	\item Comme le prouve le cas de la distance de $à$ à $\R_+\st$, il n'existe pas toujours d'élément $a\in A$ tel que $d(x,A)=d(x,a)$. Lorsqu'un tel élément existe, on dit que la distance de $x$ à $A$ est atteinte.
\end{itemize}

\begin{exo}
	Soit $B$ une boule de centre $a$ et de rayon $r>0$ d'un espace vectoriel normé $E$, et $x$ un élément de $E$. Exprimer, en fonction de $a$, $r$ et $x$, la distance de $x$ à $B$.
\end{exo}

\begin{rmq}
	Il est clair, d'après la définition, que la notion de distance à une partie dépend de la distance (ou de la norme) utilisée. L'exercice suivant permet de bien s'en rendre compte.
\end{rmq}

\begin{exo}
	Dans $\R^2$, soit $x$ le point $(-1,1)$ et $A$ la droite d'équation $y=2x$. pour chacune des normes un, deux et infinie, déterminer la distance de $x$ à $A$, et donner les points où cette distance est atteinte.
\end{exo}

\begin{exo}
	Étant données deux parties $A$ et $B$ de $E$ non vides, on pose : $$d(A,B)=\inf\{d(a,b)\ \mid\ (a,b)\in A\times B\}$$ Quels sont les axiomes d'une distance qui sont vérifiés par une telle application ?
\end{exo}

\subsection{Partie bornée, application bornée}

\subsubsection*{Partie bornée}

\begin{dfn}
	Le \textbf{diamètre} d'une partie $A$ de $E$ est défini dans $\overline{\R}_+$ par : $$\text{diam}(A)=\sup\{d(x,y)\ \mid\ (x,y)\in A^2\}$$ si $A$ est non vide, et vaut $0$ sinon.\\
	On dit qu'une partie $A$ est \textbf{bornée} si son diamètre est fini.
\end{dfn}

\paragraph{Propriétés}
\begin{enumerate}
	\item Le diamètre est une fonction croissante pour l'inclusion : $$A\subset B\implies\text{diam}(A)\leqslant\text{diam}(B)$$
	(ce qui justifie $\text{diam}(\emptyset)=0$, s'il en était besoin...).
	\item Pour $A$ et $B$ non vides, on a $\text{diam}(A\cup B)\leqslant\text{diam}(A)+\text{diam}(B)+d(A,B)$.
	\item La réunion de deux (et donc par récurrence d'un nombre fini de) parties bornées est bornée.
\end{enumerate}

\begin{ex}
	\begin{enumerate}
		\item Pour la distance discrète, toute partie est bornée.
		\item Tout segment d'un espace vectoriel normé est borné. Quel est son diamètre ?
	\end{enumerate}
\end{ex}

\begin{exo}
	\begin{enumerate}
		\item Montrer que toute boule de rayon $r>0$ d'un espace métrique est bornée de diamètre inférieur à $2r$.
		\item Donner un exemple de boule de rayon $r>0$ de diamètre strictement inférieur à $2r$.
		\item Montrer que dans un espace vectoriel normé non réduit à $\{0\}$, les boules de rayon $r>0$ sont toutes de diamètre $2r$.
	\end{enumerate}
\end{exo}

\begin{thm}{Caractérisation des parties bornées}%prop
	Soit $A$ une partie de $E$ et $a\in E$. Les propriétés suivantes sont équivalentes :
	\begin{enumerate}[label=(\roman*)]
		\item $A$ est bornée ;
		\item il existe un réel $r>0$ tel que $A\subset B_f(a,r)$.
	\end{enumerate}
\end{thm}

\begin{ex}
	On retrouve que la réunion d'un nombre fini de parties bornées est une partie bornée.
\end{ex}

\begin{thm}{}%cor
	Soit $E$ un espace vectoriel normé. Une partie $A$ de $E$ est bornée si, et seulement si : $$\exists R\in\R_+\ : \forall x\in A,\ \lVert x\rVert\leqslant R$$
\end{thm}

\subsubsection*{Application bornée, suite bornée}

\begin{dfn}
	On dit qu'une application $f$ d'un ensemble $X$ dans un espace métrique $E$ est \textbf{bornée} si $f(X)$ est une partie bornée de $E$. Dans le cas d'un espace vectoriel normé, cela revient à dire qu'il existe $R\in\R_+$ tel que : $$\forall x\in X,\ \lVert f(x)\rVert\leqslant R$$
\end{dfn}

\begin{rmq}
	\begin{itemize}
		\item Une suite d'éléments $(u_n)_{n\in\N}$ d'éléments de $E$ n'étant rien d'autre qu'une application de $\N$ dans $E$, la définition précédente s'applique.
		\item Si $E$ est un espace vectoriel normé, il est facile de vérifier que l'ensemble des applications bornées de $X$ dans $E$ est un sous-espace vectoriel du $\K$-espace vectoriel $\mathcal{F}(X,E)$. des fonctions de $X$ dans $E$.
		\item Il est clair, d'après la définition, que la notion de partie bornée (et donc de fonction ou suite bornée) dépend de la distance ou de la norme utilisée. Nous verrons :
		\begin{itemize}[label=¤]
			\item qu'il est possible qu'une partie soit bornée pour une norme mais ne le soit pas pour une autre ;
			\item que, néanmoins, deux normes équivalentes définissent les mêmes parties bornées ;
			\item en particulier, en dimension finie, on peut choisir la norme que l'on veut pour vérifier qu'une partie est bornée.
		\end{itemize}
	\end{itemize}
\end{rmq}

\section{Suites d'éléments d'un espace vectoriel normé}

On rappelle qu'une suite $a$ à valeurs dans un ensemble $X$ est une application de $\N$ dans $X$. Une telle suite est souvent notée $(a_n)_{n\in\N}$, voire plus simplement $(a_n)$. \\
L'élément $a_n$, quant à lui, est appelé \textbf{terme général} de la suite $a$.

\subsection{Suite convergente}

\begin{dfn}
	Soit $(a_n)_{n\in\N}$ une suite d'éléments de $E$. On dit que la suite $(a_n)_{n\in\N}$ converge vers $l\in E$ si la suite réelle $(d(a_n,l))_{n\in\N}$ converge vers $0$.
\end{dfn}

La suite $(a_n)_{n\in\N}$ converge donc vers $l\in E$ si, et seulement si : $$\forall\varepsilon>0,\ \exists n_0\in\N\ : \forall n\in\N,\ n\geqslant n_0\implies d(a_n,l)\leqslant\varepsilon$$
ou encore :
$$\forall\varepsilon>0,\ \exists n_0\in\N\ : \forall n\in\N,\ n\geqslant n_0\implies a_n\in B_f(l,\varepsilon)$$
Dans le cadre d'un espace vectoriel normé, cela s'écrit aussi : $$\forall\varepsilon>0,\ \exists n_0\in\N\ : \forall n\in\N,\ n\geqslant n_0\implies \lVert a_n-l\rVert\leqslant\varepsilon$$


\paragraph{Notation} Pour signifier qu'une suite $(a_n)_{n\in\N}$ tend vers $l$, on écrit $a_n\xrightarrow[n\to+\infty]{}l$ ou plus simplement $a_n\to l$.

\begin{thm}{Unicité de la limite}%prop
	Soit $(a_n)_{n\in\N}$ une suite d'éléments de $E$ ainsi que $l_1$ et $l_2$ deux éléments de $E$. Si $a_n\to l_1$ et $a_n\to l_2$, alors $l_1=l_2$.
\end{thm}
\begin{proof}
	Supposons que $(a_n)$ tende vers $l_1$ et vers $l_2$. Alors : $$\forall n\in\N,\ 0\leqslant d(l_1,l_2)\leqslant d(a_n,l_1)+d(a_n,l_2)$$ En passant, à la limite, $d(l_1,l_2)=0$ et on conclut par séparation de la distance $d$.
\end{proof}

\begin{dfn}
	Une suite $(a_n)_{n\in\N}$ est dite \textbf{convergente} s'il existe un élément $l\in E$ tel que $a_n\to l$. Cet élément $l$ est alors appelé \textbf{limite} de la suite $(a_n)_{n\in\N}$, et on le note $\lim a_n$. \\
	Une suite qui n'est pas convergente est dite \textbf{divergente}.
\end{dfn}

\paragraph{Notation} Pour dire qu'une suite est convergente (respectivement divergente), on dit souvent simplement qu'elle converge (respectivement diverge).

\paragraph{Attention} La convergence d'une suite à valeurs dans un espace métrique dépend par définition de la distance avec laquelle on travaille. Une suite peut en effet converger pour une distance et diverger pour une autre.\\
En cas d'ambiguïté, il est donc important de préciser la distance (ou norme) utilisée.

\begin{ex}
	Dans l'espace $\mathcal{C}([0,1],\R)$, on considère la suite de fonctions $(f_n)_{n\in\N}$ définie par $f_n(x)=x^n$.
	\begin{enumerate}
		\item La suite $(f_n)$ converge vers la fonction nulle au sens de la norme un mais pas au sens de la norme infinie.
		\item Supposons que $(f_n)$ converge au sens de la norme infinie et notons $f$ sa limite. Alors, $(f_n)$ converge vers $f$ au sens de la norme un, d'où une contradiction.
	\end{enumerate}
\end{ex}

\begin{exo}
	Montrer que la convergence vers $l$ dans un espace métrique $(E,d)$ est équivalente à la convergence vers $l$ dans l'espace métrique $(E,\delta)$ où $\delta=\min(1,d)$ est une distance.
\end{exo}

Exercice Moulin.EVN.37

On retiendra de cet exercice que la convergence et la limite d'une suite dépendent essentiellement de la norme choisie.

\begin{rmq}
	Nous verrons cependant que deux normes équivalentes définissent les mêmes suites convergentes. En particulier, en dimension fini, la convergence et la limite éventuelle d'une suite ne dépendant pas de la norme choisie.
\end{rmq}

\subsubsection*{Opérations sur les suites convergentes dans un espace vectoriel normé}

Ici, $E$ est un espace vectoriel normé.

\begin{thm}{Combinaison linéaire de deux suites convergentes}%prop
	Si deux suites $(a_n)$ et $(b_n)$ convergent respectivement vers $l_1$ et $l_2$, et si $\lambda$ et $\mu$ sont deux scalaires, alors la suite de terme général $\lambda a_n+\mu b_n$ converge vers $\lambda l_1+\mu l_2$.
\end{thm}

\begin{rmq}
	Par récurrence, on généralise le résultat précédente à une combinaison linéaire d'un nombre fini de suites : si l'on dispose de $p$ suites, de termes généraux respectifs $a_n^{(1)},\ldots,a_n^{(p)}$, convergeant respectivement vers $l_1;\ldots,l_p$, et si $\lambda_1,\ldots,\lambda_p$ sont des scalaires, alors la suite de terme général $\displaystyle\sum_{k=1}^{p}\lambda_ka_n^{(k)}$ converge vers $\displaystyle\sum_{k=1}^{p}\lambda_k l_k$.
\end{rmq}

\begin{thm}{}%cor
	L'ensemble des suites convergentes à valeurs dans $E$ est un sous-espace vectoriel de l'espace $E^\N$ des suites à valeurs dans $E$ sur lequel l'application $u\mapsto\lim u$ est linéaire.
\end{thm}

\begin{thm}{}%prop
	Si une suite $(a_n)_{n\in\N}$ converge vers $l$, alors la suite de terme général $\lVert a_n\rVert$ converge vers $\lVert l\rVert$.
\end{thm}

\begin{thm}{}%cor
	Toute suite convergente est bornée.
\end{thm}

\begin{thm}{}%prop
	Soit $(u_n)_{n\in\N}$ une suite de $E$ et $(\lambda_n)_{n\in\N}$ une suite de $\K$, convergeant respectivement vers $a\in E$ et $\lambda\in\K$. Alors, $(\lambda_n u_n)_{n\in\N}$ converge vers $\lambda a$.
\end{thm}

\begin{rmq}
	On verra ultérieurement une généralisation de ce résultat à d'autres types de produits.
\end{rmq}

\subsection{Suite à valeurs dans un espace produit}

Soit $(E_1,\lVert\cdot\rVert_1),\ldots,(E_pn\lVert\cdot\rVert_p)$ des espaces vectoriels normés. On rappelle que l'espace $E_1\times\cdots\times E_p$ est muni d'une structure d'espace vectoriel normé grâce à la norme produit définie par : $$\lVert(x_1,\ldots,x_p)\rVert=\max(\lVert x_1\rVert_2,\ldots,\lVert x_p\rVert_p).$$ Pour tout $k\in\llbracket1,p\rrbracket$, on considère $(a_n^{(k)})_{n\in\N}$ une suite à valeurs dans $E_k$, et on s'intéresse à la convergence au sens de la norme produit de la suite $(a_n)_{n\in\N}$ à valeurs dans $E_1\times\cdots\times E_p$ de terme général : $$a_n=(a_n^{(1)},\ldots,a_n^{(p)}).$$ La proposition suivante énonce que la convergence de la suite $(a_n)_{n\in\N}$ revient à celle de chacune des suites $(a_n^{(k)})_{n\in\N}$.

\begin{thm}{}%prop
	La suite de terme général $(a_n^{(1)},\ldots,a_n^{(p)})$ converge si, et seulement si, pour tout $k\in\llbracket1,p\rrbracket$, la suite $(a_n^{(k)})_{n\in\N}$ converge.\\
	De plus, en notant alors $l_k$ la limite de la suite $(a_n^{(k)})_{n\in\N}$, la suite $(a_n)_{n\in\N}$ converge vers l'élément $(l_1,\ldots,l_p)$ de $E_1\times\cdots\times E_p$.
\end{thm}

\begin{ex}
	Une suite à valeurs dans $\K^p$ converge au sens de la norme infinie si, et seulement si, chacune de ses composantes converge. L'étude de la convergence d'une telle suite revient donc à l'étude de $p$ suites à valeurs dans $\K$.
\end{ex}

\begin{rmq}
	Le fait que toutes les normes soient équivalentes en dimension finie entraîne que l'étude de la convergence, au sens de n'importe quelle norme (par exemple la norme un ou la norme deux) d'une suite à valeurs dans $\K^p$ revient à l'étude des suites composantes.
\end{rmq}

\paragraph{Extension} Tout ce qui vient d'être fait dans un produit (fini) d'espaces vectoriels normés peut sans problème se généraliser dans un produit (fini) d'espaces métriques.

\subsection{Suite extraite, valeur d'adhérence}

Comme dans le cas des suites à valeurs réelles ou complexes, on appelle \textbf{suite extraite} ou \textbf{sous-suite} d'une suite $(a_n)_{n\in\N}$ toute suite de la forme $(a_{\varphi(n)})_{n\in\N}$ où $\varphi$ est une extraction, c'est-à-dire une application strictement croissante de $\N$ dans $\N$.

\begin{thm}{}%prop
	Soit $u$ une suite de $E$ convergeant vers $l\in E$. Alors toute sous-suite de $u$ converge également vers $l$.
\end{thm}

Exercice Moulin.EVN.38

\begin{dfn}
	On appelle \textbf{valeur d'adhérence} d'une suite $(a_n)_{n\in\N}$ tout élément de $E$ qui est limite d'une sous-suite convergente de $(a_n)_{n\in\N}$.
\end{dfn}

\begin{rmq}
	D'après la proposition précédente, une suite convergente possède une unique valeur d'adhérence.
\end{rmq}

\begin{met}
	Pour montrer qu'une suite est divergente, on peut utiliser la contraposée de la remarque précédente : il suffit de montrer que la suite considérée possède au moins deux valeurs d'adhérence (ou n'en possède aucune).
\end{met}

\begin{ex}
	La suite de terme général $e^{in}$ est divergente, car si $l$ est une valeur d'adhérence, alors $le^{i}$ aussi...
\end{ex}

\begin{rmq}
	Comme pour les suites réelles, la convergence des deux suites $(u_{2n})$ et $(u_{2n+1})$ vers la même limite $l$ entraîne la convergence de $(u_n)$ vers $l$.
\end{rmq}

\begin{thm}{Caractérisation des valeurs d'adhérence}%prop
	Soit $(a_n)_{n\in\N}$ une suite à valeurs dans $E$ et $x$ un élément de $E$. Alors $x$ est valeur d'adhérence de la suite $(a_n)_{n\in\N}$ si, et seulement si : $$\forall\varepsilon>0,\ \forall n\in\N,\ \exists p\geqslant n\ : \ d(a_n,x)\leqslant\varepsilon$$
\end{thm}

\paragraph{Traduction en termes de boules} Soit $u=(u_n)$ une suite de $E$ et $l$ un élément de $E$.
\begin{itemize}
	\item La convergence de $u$ vers $l$ se traduit par le fait que toute boulée centrée en $l$ et de rayon strictement positif contient tous les termes de $u$ à partir d'un certain rang.
	\item L'élément $l$ est valeur d'adhérence de $u$ si, et seulement si, $u$ prend une infinité de valeurs dans toute boule centrée en $l$ et de rayon strictement positif.
\end{itemize}

\section{Séries à valeurs dans un espace vectoriel normé}

Ici, $E$ désigne un espace vectoriel normé.

\subsection{Suites et séries}

A toute suite $(u_n)_{n\in\N}\in\E^\N$, on associe la suite $(S_n)_{n\in\N}$ définie par : $$\forall n\in\N,\ S_n=\sum_{k=0}^{n}u_k$$ L'application $u\mapsto S$ est un isomorphisme d'espace vectoriel de $\E^\N$. Sa réciproque est l'application $S\mapsto u$, où $u$ est définie par : $$u_0=S_0\ \ \text{et}\ \ \forall n\in\N\st,\  u_n=S_n-S_{n-1}$$ Pour chaque élément de $\E^\N$, on peut donc considérer au choix la suite $(u_n)_{n\in\N}$ ou la suite $(S_n)_{n\in\N}$ associée.

\begin{dfn}
	Si $u\in\E^\N$, et $S_n=\displaystyle\sum_{p=0}^{n}u_p$, on appelle série de terme général CENSURE. \\
	On la note $\displaystyle\sum u_n$. \\
	Le \textbf{terme général} est donc le vecteur $u_n$ ; les \textbf{sommes partielles} sont les $S_n$.
\end{dfn}

\begin{dfn}
	On dit que la série $\displaystyle\sum u_n$ \textbf{converge} si la suite $(S_n)$ admet une limite dans $\E$, appelée \textbf{somme de la série} et notée $\displaystyle\sum_{n=0}^{+\infty}u_n$. Dans le cas contraire, on dit que la série \textbf{diverge}. On parle de la nature d'une série pour parler de sa convergence ou sa divergence.
\end{dfn}

\paragraph{Notation} On note $S_n(u)$ la $n$-ième somme partielle de la série de terme général $u_n$. Lorsque la série converge, on note aussi $S(u)$ sa somme.

\begin{rmq}
	\begin{itemize}
		\item Lorsque la suite $u$ est définie seulement à partir d'un certain rang $n_0$, la série de terme général $u_n$ est notée $\displaystyle\sum_{n\geqslant n_0}u_n$ ; il s'agit de la suite des sommes partielles définie par : $$\forall n\geqslant n_0,\ S_n=\sum_{p=n_0}^{n}u_p$$
		\item Le caractère asymptotique de la notion de limite permet de voir que pour $n_0\in\N$, une série $\displaystyle\sum u_n$ converge si, et seulement si, la série $\displaystyle\sum_{n\geqslant n_0}u_n$.
	\end{itemize}
\end{rmq}

\begin{met}
	Pour étudier la convergence d'une suite $(u_n)_{n\in\N}$, il faut penser à étudier éventuellement la convergence de la série $\displaystyle\sum_{n\geqslant 1}(u_{n+1}-u_n)$.
\end{met}

\begin{thm}{}%prop
	Si une série converge, alors son terme général converge vers $0$.
\end{thm}

\begin{dfn}
	Une série dont le terme général ne tend pas vers $à$ est dite \textbf{grossièrement divergente}.
\end{dfn}

\paragraph{Attention} Comme pour les séries numériques, la convergence vers $0$ du terme général est une condition nécessaire mais non suffisante de convergence.

\begin{thm}{}%prop
	L'ensemble des suites $(u_n)_{n\in\N}\in\E^\N$ telles que la série $\displaystyle\sum u_n$ converge est un sous-espace vectoriel de $\E^\N$ sur lequel l'application somme est linéaire.
\end{thm}

\begin{dfn}
	Lorsque la série $\displaystyle\sum u_n$, on note $\displaystyle R_n(u)=\sum_{p=n+1}^{+\infty}u_p$, qu'on appelle \textbf{reste} au rang $p$.
\end{dfn}

\begin{rmq}
	Dans ce cas, la suite $(R_n(u))$ converge vers $0$ comme le prouve la relation : $$R_n(u)=S(u)-S_n(u)$$
\end{rmq}

\subsection{Convergence absolue}

\begin{dfn}
	On dit que la série $\displaystyle\sum u_n$ est \textbf{absolument convergente}, ou \textbf{converge absolument}, si la série réelle positive $\displaystyle\sum \lVert u_n\rVert$ est convergente.
\end{dfn}

\paragraph{Attention} Contrairement à ce qui se passe dans $\R$ ou dans $\C$, une série absolument convergente peut très bien ne pas converger.

\begin{ex}
	Dans $\K[X]$ muni de la norme $\displaystyle\left\lVert\sum_{k=0}^{+\infty}a_kX^k\right\rVert=\max\lvert a_k\rvert$, la série $\displaystyle\sum\frac{X^n}{n!}$ converge absolument, mais est divergente.
\end{ex}

Néanmoins, nous démontrons ultérieurement le résultat suivant :

\begin{thm}{}%thm
	Toute série absolument convergente d'un espace vectoriel normé de dimension finie est convergente.
\end{thm}

\section{Topologie d'un espace métrique}

\subsection{Parties ouvertes}

\begin{dfn}
	Soit $U$ une partie de $E$. On dit que $U$ est un \textbf{ouvert} de $E$, ou une \textbf{partie ouverte} de $E$, si l'une des deux assertions équivalentes suivantes est vérifiée :
	\begin{enumerate}[label=(\roman*)]
		\item pour tout $x\in U$, il existe $r>0$ tel que la boule ouverte $B_o(x,r)$ soit incluse dans $U$ ;
		\item pour tout $x\in U$, il existe $r>0$ tel que la boule fermée $B_f(x,r)$ soit incluse dans $U$.
	\end{enumerate}
\end{dfn}

L'équivalence entre les deux assertions de la définition ci-dessus est évidente :
\begin{itemize}
	\item si la boule ouverte centrée en $x$ et de rayon $r$ est incluse dans $U$, alors il en est de même pour la boule fermée centrée en $x$ et de rayon $r/2$ ;
	\item si la boule fermée centrée en $x$ et de rayon $r$ est incluse dans $U$, alors il en est de même pour la boule ouverte centrée en $x$ et de rayon $r$.
\end{itemize}

\begin{ex}
	Il est clair, d'après la définition, que $E$ et $\emptyset$ sont des parties ouvertes de $E$.
\end{ex}

\begin{thm}{}%prop
	Toute boule ouverte est une partie ouverte de $E$.
\end{thm}

\begin{proof}
	Pour $x$ appartenant à une boule ouvert de centre $a$ et de rayon $R$, considérer la boule ouverte de centre $x$ et de rayon $R-d(x,a)$. Faire un dessin peu aider...
\end{proof}

\begin{rmq}
	Le dessin fait lors de la démonstration précédente correspond à une norme euclidienne, mais nous a permis de trouver une démonstration valable pour tout norme. Cette représentation des boules n'induit en général pas de fausse intuition lors de considérations topologiques, et est donc couramment utilisée. Concrètement, le risque majeur de cette vision euclidienne concerne le cas d'égalité dans l'inégalité triangulaire.
\end{rmq}

\paragraph{Attention} Le caractère ouvert ou non d'un ensemble n'est pas une notion intrinsèque à cet ensemble, mais dépend de l'espace dans lequel on le considère. Si $F$ est un sous-espace métrique de $E$, alors une partie ouverte ouvert dans $F$ n'est en général pas une partie ouvert dans $E$, comme le montre l'exercice suivant.

\begin{exo}
	Soit $a$ et $b$ deux réels tels que $a<b$.
	\begin{enumerate}
		\item L'intervalle ouvert $]a,b[$ est-il un ouvert de $\R$ (muni de la valeur absolue) ?
		\item L'intervalle ouvert $]a,b[$ est-il un ouvert de $\C$ (muni du module) ?
	\end{enumerate}
\end{exo}

\begin{rmq}
	Plus généralement, tout intervalle ouvert non vide de $\R$, c'est-à-dire de la forme $]-\infty,b[$, $]a,b[$ ou $]a,+\infty[$ est un ouvert de $\R$, mais pas de $\C$.
\end{rmq}

\begin{thm}{}%prop
	\begin{itemize}
		\item La réunion d'une famille quelconque d'ouverts de $E$ est un ouvert de $E$.
		\item L'intersection d'une famille finie d'ouverts de $E$ est un ouvert de $E$.
	\end{itemize}
\end{thm}

\begin{rmq}
	En particulier, une réunion quelconque de boules ouvertes est un ouvert de $E$. Réciproquement :
\end{rmq}

\begin{exo}
	Montrer que toute partie ouvert peut s'écrire comme une réunion de boules ouvertes.
\end{exo}

\paragraph{Attention} Avant d'affirmer qu'une intersection d'ouverts est un ouvert, il faut bien vérifier le caractère fini de cette intersection. En effet, une intersection quelconque de parties ouvertes n'est en général pas ouverte, comme l'illustre l'exemple suivant.

\begin{exo}
	Pour tout $n\in\N\st$, on note $U_n=\displaystyle\left]-\frac{1}{n},\frac{1}{n}\right[$. Bien que chacun des $U_n$ soit un ouvert, leur intersection $\{0\}$ ne l'est pas.
\end{exo}

\begin{thm}{}%prop
	Soit $E_1,\ldots,E_p$ des espaces métriques. On munit l'espace produit $E_1\times\cdots\times E_p$ de la distance produit. Si $U_1,\ldots,U_p$ sont ouverts respectivement dans $E_1,\ldots,E_p$, alors $U_1\times\cdots\times U_p$ est ouvert dans $E_1\times\cdots\times E_p$.
\end{thm}

\subsection{Parties fermées}

\begin{dfn}
	On dit qu'une partie de $E$ est un \textbf{fermé} de $E$, ou une \textbf{partie fermée} de $E$, si son complémentaire dans $E$ est ouvert.
\end{dfn}

\begin{ex}
	D'après la définition, il est clair que $E$ et $\emptyset$ sont des parties fermées de $E$, puisque leurs complémentaires sont des ouverts de $E$.
\end{ex}

\begin{exo}
	Montrer que tout intervalle fermé de $\R$, c'est-à-dire tout intervalle de la forme $]-\infty,b]$, $[a,b]$ ou $[a,+\infty[$, est un fermé de $\R$.
\end{exo}

\paragraph{Attention} Il ne faut pas croire, par jeu de mots, qu'une partie qui n'est pas ouvert est nécessairement fermée. En effet, dans tout $\K$-espace vectoriel normé non réduit à $\{0\}$, il existe des parties ni ouvertes ni fermées.\\
C'est le cas par exemple de l'intervalle $[0,1[$ dans $\R$.\\
De plus, il existe des parties à la fois ouvertes et fermées, comme $E$ et $\emptyset$.

\begin{exo}
	Quelles sont les parties à la fois ouvertes et fermées d'un ensemble muni de la distance discrète ?
\end{exo}

\begin{thm}
	Toute boule fermée est une partie fermée de $E$.
\end{thm}
\begin{proof}
	On montre que le complémentaire de toute boule fermée est un ouvert. Pour $x$ n'appartenant pas à la boule fermée de centre $a$ et de rayon $R$, considérer la boule ouvert centrée en $x$ et de rayon $d(x,a)-R$. Faire un dessin peut aider...
\end{proof}

\begin{ex}
	\begin{enumerate}
		\item Tout singleton est un fermé de $E$.
		\item Tout segment  de $\R$ est un fermé de $\R$. En effet, on peut l'interpréter comme une boule fermée.
	\end{enumerate}
\end{ex}

\begin{thm}{}%prop
	\begin{itemize}
		\item L'intersection d'une famille quelconque de fermés est un fermé.
		\item La réunion d'une famille finie de fermés est un fermé.
	\end{itemize}
\end{thm}

\paragraph{Attention} Avant d'affirmer qu'une réunion de fermés est fermée, il faut bien vérifier le caractère fini de cette réunion. En effet, une réunion quelconque de fermés n'est en général pas fermée, comme l'illustre l'exercice suivant.

\begin{exo}
	On se place dans $\R$. Pour tout $n\in\N\st$, on note $F_n=\left[\frac{1}{n},1-\frac{1}{n}\right]$. Déterminer l'ensemble $A=\displaystyle\bigcup_{n\in\N\st}F_n$. Cet ensemble est-il fermé ?
\end{exo}

\begin{rmq}
	Plus généralement, toute partie $A$ peut-être obtenue comme une réunion de fermés : il suffit d'écrire $A=\displaystyle\bigcup_{a\in A}\{a\}$.
\end{rmq}

\begin{ex}
	Montrer, en l'obtenant comme une intersection de fermés, que toute sphère est une partie fermée.
\end{ex}
%boule fermée intersectée avec le complémentaire de la boule ouverte

\subsection{Voisinage d'un point}

\begin{dfn}
	Étant donnée une partie $A$ de $E$ et $a\in E$, on dit que $A$ est un \textbf{voisinage de $a$} s'il existe $r>0$ tel que la boule ouvert $B_o(a,r)$ soit incluse dans $A$. \\
	On note $\mathcal{V}_a$ l'ensemble des voisinages de $a$.
\end{dfn}

\begin{rmq}
	De manière évidente, pour que $A$ soit un voisinage de $a$, il est nécessaire que $a$ appartienne à $A$.
\end{rmq}

Dire que $A$ est un voisinage de $a$ signifie qu'il existe $r>0$ tel que :
$$\forall x\in E,\ d(x,a)<r\implies x\in A,$$ autrement it, se placer suffisamment proche du point $a$ nous assure d'être dans $A$. \\
En revanche, il peut y avoir dans les voisinages de $a$ des points "très loin" de $a$.

\begin{rmq}
	On peut, sans en changer le sens, remplacer "boule ouvert" par "boule fermée" dans la définition qui précède.
\end{rmq}

\begin{ex}
	Une partie de $E$ est un ouvert de $E$ si, et seulement si, elle est voisinage de tous ses points.
\end{ex}

\begin{exo}
	Montrer que les ouverts, les fermés et les voisinages d'un espace métrique $(E,d)$ sont les mêmes que leurs équivalents dans $(E,\delta)$ où $\delta=\min(1,d)$. \\
	On dit que les espaces métriques $(E,d)$ et $(E,\delta)$ ont la \textbf{même topologie}.
\end{exo}

\subsubsection*{Opérations sur les voisinages}

\begin{thm}{}%prop
	Soit $a$ un élément de $E$.
	\begin{itemize}
		\item Une intersection finie de voisinages de $a$ est encore un voisinage de $a$.
		\item Si $A$ est un voisinage de $A$, alors toute partie contenant $A$ est également un voisinage de $A$.
	\end{itemize}
\end{thm}

\begin{rmq}
	Le deuxième point de la proposition précédente permet d'affirmer que la réunion d'une famille non vide quelconque de voisinages de $a$ est un voisinage de $a$.
\end{rmq}

\begin{exo}
	Montrer qu'une partie $A$ est un voisinage d'un point $a$ si, et seulement s'il existe un ouvert contenant $a$ et inclus dans $A$.
\end{exo}

\subsubsection*{Limites et voisinages}

\begin{thm}{}%prop
	Soit $a\in E$. Une suite $(u_n)$ de $E$ converge vers $a$ si, et seulement si : $$\forall V\in\mathcal{V}_a,\ \exists n_0\in\N\ : \ n\geqslant n_0\implies u_n\in V$$
\end{thm}

\begin{exo}
	Traduire de même le fait que $a$ soit une valeur d'adhérence de $(u_n)$. Comparer avec les traductions en termes de boules. Explication ?
\end{exo}

\subsection{Extension (HP) aux voisinages de l'infini}

On suppose que $E$ est ici un espace vectoriel normé.

\begin{dfn}
	\begin{enumerate}
		\item On dit que $V\subset E$ est un \textbf{voisinage de l'infini} s'il existe $A$ tel que : $$\lVert x\rVert\geqslant A\implies x\in V$$
		On note $\mathcal{V}_\infty$.
		\item On dit que $V\subset\R$ est un \textbf{voisinage de $+\infty$} s'il existe $A$ tel que $[A,+\infty[\subset V$. \\
		On note $\mathcal{V}_{+\infty}$ l'ensemble des voisinages de $+\infty$.
		\item On dit que $V\subset\R$ est un \textbf{voisinage de $-\infty$} s'il existe $A$ tel que $]-\infty,A]\subset V$. \\
		On note $\mathcal{V}_{-\infty}$ l'ensemble des voisinages de $-\infty$.
	\end{enumerate}
\end{dfn}

\begin{rmq}
	\begin{itemize}
		\item Les voisinages de l'infini, de $+\infty$ et de $-\infty$ vérifient les opérations sur les voisinages énoncées dans la section précédente.
		\item Dans $\R$, a-t-on $\mathcal{V}_\infty=\mathcal{V}_{+\infty}\cap\mathcal{V}_{-\infty}$ $\mathcal{V}_\infty=\mathcal{V}_{+\infty}\cup\mathcal{V}_{-\infty}$ ?	\end{itemize}
\end{rmq}

\begin{dfn}
	Une suite $(u_n)_{n\in\N}$ d'un espace vectoriel normé $E$ \textbf{tend vers l'infini}, ou \textbf{tend vers $\infty$} (sans signe), et l'on écrit $\lim u_n=\infty$, lorsque la suite réelle $(\lVert u_n\rVert)_{n\in\N}$ tend vers $+\infty$.
\end{dfn}

\begin{thm}{Écriture topologique des limites}%prop
	Une suite $(u_n)_{n\in\N}$ tend vers $l$ si, et seulement si, l'on a : $$\forall V\in\mathcal{V}_l,\ \exists n_0\in\N\ : \ n\geqslant n_0\implies u_n\in V,$$ ceci étant valable pour $l\in E$ ou $l=\infty$ lorsque $l$ est un espace vectoriel normé, et pour $l\in\R$ ou $l=\pm\infty$ lorsque $E=\R$.
\end{thm}

\begin{rmq}
	Cette proposition permet d'unifier les définitions et de prouver avec moins d'effort le théorème de composition des limites.
\end{rmq}

\subsection{Point intérieur, intérieur d'une partie}

Désormais, $E$ est de nouveau un espace métrique.

\begin{dfn}
	On dit qu'un point $x$ de $E$ est intérieur à une partie $A$ de $E$ si $A$ est un voisinage de $x$.
\end{dfn}

\begin{dfn}
	L'\textbf{intérieur de $A$} est le plus grand ouvert inclus dans $A$. On le note $\mathring{A}$ ou $\text{Int}(A)$.
\end{dfn}

\begin{thm}
	L'intérieur de $A$ est l'ensemble de ses points intérieurs.
\end{thm}

On rappelle, d'après la définition des voisinages, que dire que $x$ est intérieur à $A$ signifie qu'on peut trouver $r>0$ tel que la boule centrée en $x$ et de rayon $r$ soit incluse dans $A$.

\paragraph{Propriétés} Les propriétés suivantes peuvent s'obtenir à l'aide de l'une ou l'autre des deux définitions de l'intérieur et des points intérieurs.
\begin{itemize}
	\item Si $A$ et $B$ sont deux parties vérifiant $A\subset B$, alors on a $\mathring{A}\subset\mathring{B}$.
	\item L'intérieur d'une réunion contient la réunion des intérieurs.
	\item L'intérieur d'une intersection finie est contenue dans l'intersection des intérieurs, avec égalité pour une intersection finie.
\end{itemize}

\begin{exo}
	Montrer qu'une partie $A$ est ouvert si, et seulement si,  $\mathring{A}=A$.
\end{exo}

\subsection{Point adhérent, adhérence d'une partie}

\begin{dfn}
	On dit qu'un point $x\in E$ est \textbf{adhérent} à une partie $A$ de $E$ si pour tout $r>0$, on a $B_o(x,r)\cap A\ne\emptyset$.
\end{dfn}

\begin{ex}
	\begin{enumerate}
		\item Les points adhérents à $A\ne\emptyset$ sont les points à distance nulle de $A$.
		\item Borne supérieur d'une partie de $\R$.
	\end{enumerate}
\end{ex}

\begin{thm}{}%prop
	un point $a$ est adhérent à $A$ si, et seulement si tout voisinage de $a$ rencontre $A$, ce qui s'écrit : $$\forall V\in\mathcal{V}_a,\ V\cap A\ne\emptyset$$
\end{thm}

\begin{dfn}
	L'adhérence de $A$ est le plus petit fermé de $E$ contenant $A$. On le note $\overline{A}$ ou $\text{Adh}(A)$.
\end{dfn}

\begin{ex}
	Une partie de $A$ est un fermé si, et seulement si, $\overline{A}=A$.
\end{ex}

\begin{exo}
	Montrer que l'ensemble des valeurs d'adhérence d'une suite est un fermé.
\end{exo}

Exercice Moulin.EVN.33

\begin{thm}{}%prop
	L'adhérence de $A$ est l'ensemble de ses points adhérents.
\end{thm}

\begin{rmq}
	Dire qu'un point $x$ n'appartient pas à l'adhérence de $A$ signifie qu'on peut trouver $r>0$ tel que $B_o(x,r)\cap A=\emptyset$ (ou encore $B_o(x,r)\subset E\setminus A$), ce qui signifie que $x$ appartient à l'intérieur de $E\setminus A$. On a donc la relation : $$E\setminus\left(\text{Adh}(A)\right)\text{Int}(E\setminus A),$$
	ce qui mène également aux deux relations suivantes : $$\text{Adh}(A)=E\setminus\text{Int}(E\setminus A)\ \ \text{et}\ \ \text{Int}(A)=E\setminus\text{Adh}(E\setminus A)$$
\end{rmq}

\subsubsection*{Caractérisation séquentielle des points adhérents}

Le résultat suivant donne une caractérisation très importante des points adhérents à une partie $A$ d'un espace \textbf{métrique.}

\begin{thm}{Caractérisation séquentielle des points adhérents}
	Un point $x$ est adhérent à une partie $A$ si, et seulement s'il existe une suite d'éléments de $A$ qui converge vers $x$.
\end{thm}

Ce résultat s'énonce ainsi : l'adhérence d'une partie $A$ d'un espace métrique est l'ensemble des limites des suites convergentes à valeurs dans $A$.

\begin{ex}
	Borne supérieure d'une partie de $\R$.
\end{ex}

Les notions de point adhérent à une partie et de valeur d'adhérence d'une suite sont des notions bien distinctes, comme le montre l'exercice suivant.

\begin{exo}
	Soit $(u_n)_{n\in\N}$ une suite à valeurs dans $E$. Pour $p\in\N$, on note $A_p$ l'ensemble $\{u_n\mid n\geqslant p\}$. Montrer que l'ensemble des valeurs d'adhérence de la suite $(u_n)_{n\in\N}$ est :
	$$\bigcap_{p\in\N}\overline{A_p}$$ Retrouver le résultat de l'exercice précédent.
\end{exo}

\subsubsection*{Caractérisation séquentielle des parties fermées}

Une partie étant un fermé si, et seulement si, elle est égale à sa propre adhérence, on déduit de la caractérisation séquentielle des points adhérents le résultat suivant :

\begin{thm}{Caractérisation séquentielle des parties fermées}
	Une partie $A$ est un fermé si, et seulement si, la limite de toute suite convergente d'éléments de $A$ appartient à $A$.
\end{thm}

\subsubsection*{Extension (HP) de la notion de point adhérent}

On étend la définition de point adhérent à une partie $A$ de $E$ dans les cas suivants :
\begin{itemize}
	\item $a\in E$ ;
	\item $a=\pm\infty$ pour $E=\R$ ;
	\item $a=\infty$ pour un espace vectoriel normé $E$ quelconque.
\end{itemize}

\begin{dfn}
	 Un tel élément est adhérent à $A$ s'il existe une suite d'éléments de $A$ qui tend vers $a$.
\end{dfn}

Ainsi :
\begin{itemize}
	\item $-\infty$ est adhérent à $A$ si, et seulement si, $A$ est une partie de $\R$ non minorée ;
	\item $+\infty$ est adhérent à $A$ si, et seulement si, $A$ est une partie de $\R$ non majorée ;
	\item $\infty$ est adhérent à $A$ si, et seulement si, $A$ est non bornée.
\end{itemize}

\begin{thm}{}%prop
	Avec les définitions précédentes, $a$ est adhérent à $A$ si, et seulement si, tout voisinage de $a$ rencontre $A$, ce qui s'écrit : 
	$$\forall V\in\mathcal{V}_a,\ V\cap A\ne\emptyset$$
\end{thm}

\subsection{Densité}

\begin{dfn}
	Soit une partie de $A$. On dit que $D$ est \textbf{dense dans $A$} si l'une des trois propriétés équivalentes suivantes est vérifiée :
	\begin{enumerate}[label=(\roman*)]
		\item l'adhérence de $D$ contient $A$ :
		\item pour tout $a\in A$ et pour tout $r>0$, il existe $x\in D$ tel que $d(x,a)\leqslant r$ ;
		\item pour tout $a\in A$, il existe une suite d'éléments qui converge vers $a$.
	\end{enumerate}
\end{dfn}
\begin{proof}
	Dans la définition précédente, les propriétés $(ii)$ et $(iii)$ ne sont que des reformulations de la propriété $A\subset\overline{D}$.
\end{proof}

\begin{ex}
	\begin{enumerate}
		\item Les ensembles $\Q$, $\R\setminus\Q$ et $\mathbb{D}$ sont denses dans $\R$.
		\item L'ensemble $\mathcal{GL}_n(\K)$ est dense dans $\mathcal{M}_n(\K)$.
		\item Soit $(a,b)\in\R^2$ tel que $a<b$. L'ensemble des fonctions en escalier est dense dans $\mathcal{CM}([a,b],\K)$ au sens de la norme infinie.
		\item Nous verrons que l'ensemble des fonctions polynomiales sur $[a,b]$ est dense dans $\mathcal{C}([a,b],\K)$ au sens de la norme infinie.
	\end{enumerate}
\end{ex}

Exercice Moulin.EVN.34

\subsection{Frontière}

\begin{dfn}
	Soit $A$ une partie de $E$. On appelle \textbf{frontière de $A$}, et on note $\text{Fr}(A)$, l'ensemble défini par : $$\text{Fr}(A)=\overline{A}\setminus\mathring{A}.$$
\end{dfn}

Dire qu'un point $x$ appartient à la frontière de $A$ signifie donc que $x\in\overline{A}$ et $x\notin\mathring{A}$, c'est-à-dire que pour tout $r>0$, on a : 
$$B_o(x,r)\cap A\ne\emptyset\ \ \text{et}\ \ B_o(x,r)\cap(E\setminus A)\ne\emptyset$$

\begin{ex}
	\begin{enumerate}
		\item Soit $E$ un espace vectoriel normé non trivial. Si $B$ désigne la boule fermée de centre $a\in E$ et de rayon $r>0$, alors l'intérieur de $B$ est la boule ouvert de mêmes centre et rayon, et la frontière de $B$ est alors la sphère de mêmes centre et rayon.
		\item Soit $a$ et $b$ deux réels tels que $a<b$. Quelle est la frontière d'un intervalle dont les bornes sont $a$ et $b$ ?
	\end{enumerate}
\end{ex}

\begin{exo}
	Quelle est la frontière de $\Q$ dans $\R$ ?
\end{exo}

\begin{exo}
	Soit $A$ une partie de $E$. Montrer : $\text{Fr}(A)=\overline{A}\setminus\overline{(E\setminus A)}$.
\end{exo}

\begin{rmq}
	Du résultat de l'exercice précédent, il résulte que la frontière de $A$ est une partie fermée comme intersection de (deux) parties fermées.
\end{rmq}

\begin{exo}
	\begin{itemize}
		\item Quelle est la frontière des parties d'un ensemble muni de la distance discrète ?
		\item Quelle sont les parties de $\R$ de frontière vide ?
	\end{itemize}
\end{exo}

\begin{thm}{}%prop
	Pour toute partie $A$ de $E$, les trois parties : $$\mathring{A},\ \text{Fr}(A)\ \text{et}\ E\setminus\overline{A}$$ forment un recouvrement disjoint de $E$.
\end{thm}

\subsection{Ouvert relatif, fermé relatif, voisinage relatif}

"Il va falloir parler de cette notion qui est assez vide."

\begin{dfn}
	Soit $A$ une partie de $E$.
	\begin{itemize}
		\item Une partie $U$ de $A$ est un \textbf{ouvert relatif} à $A$ s'il existe un ouvert $\tilde{U}$ de $E$ tel que $U=\tilde{U}\cap A$.
		\item Une partie $F$ de $A$ est un \textbf{fermé relatif} à $A$ s'il existe un fermé $\tilde{F}$ de $E$ tel que $F=\tilde{F}\cap A$.
		\item Soit $a$ un point de $A$. Une partie $V$ de $A$ est un \textbf{voisinage relatif} à $A$ du point $a$ s'il existe un voisinage $\tilde{V}$ de $a$ dans $E$ tel que $V=\tilde{V}\cap A$.
	\end{itemize}
\end{dfn}

On constate facilement que si $A$ est égal à $E$, alors ces notions coïncident avec les notions déjà vues d'ouvert, de fermé et de voisinage. En revanche, ce n'est pas le cas si $A$ est quelconque, comme illustré ci-dessous.

\begin{ex}
	L'intervalle $[0,1[$, bien qu'il ne soit pas ouvert dans $\R$, est un ouvert relatif à $[0,1]$.
\end{ex}

\begin{exo}
	On se place dans $\R$. Montrer que l'intervalle $]0,1]$, bien qu'il ne soit pas fermé dans $\R$, est un fermé relatif à $]0,2]$.
\end{exo}

\begin{exo}
	On travaille dans $\C$. Montrer que l'intervalle $[0,2]$, bien qu'il ne soit pas un voisinage du point $1$, en est un voisinage relatif à $\R$.
\end{exo}

\begin{ex}
	Soit $F$ un fermé de $E$. Les fermés relatifs à $F$ sont donc les parties fermées de $E$ contenues dans $F$. On dispose d'un résultat similaire pour les ouverts.
\end{ex}

\subsubsection*{Remarque importe pour saisir la notion (donc HP)}

En fait, les ouverts, fermés et voisinages relatifs à $A$ ne sont pas autre chose que les ouverts, fermés et voisinages de l'espace métrique $A$ muni de la distance induite par celle de $E$.

\begin{exo}
	Montrer que toute partie $A$ de $E$ est à la fois un ouvert relatif à $A$ et un fermé relatif à $A$.
\end{exo}

\subsection{Espaces topologiques (HP)}

On donne ici une présentation d'espaces plus généraux, dits topologiques, dans lesquels les notions d'ouverts, de fermés, de voisinages et de convergence sont parfaitement adaptées. Nous verrons que cette notion généralise la notion d'espace métrique.

\begin{dfn}
	Soit $E$ un ensemble. On appelle \textbf{topologie} sur $E$ toute partie $\mathcal{O}$ de $\mathcal{P}(E)$ telle que :
	\begin{enumerate}[label=(\roman*)]
		\item $\emptyset\in\mathcal{O}$ et $E\in\mathcal{O}$ ;
		\item $\mathcal{O}$ est stable par réunion quelconque ;
		\item $\mathcal{O}$ est stable par intersection finie.
	\end{enumerate}
	Un tel couple $(E,\mathcal{O})$ est appelé \textbf{espace topologique}, et alors, les éléments de $\mathcal{O}$ sont appelés les \textbf{ouverts} de $E$.
\end{dfn}

\begin{ex}
	$\mathcal{O}=\mathcal{P}(E)$ est la \textbf{topologie discrète}. Elle est métrisable par la distance discrète.
\end{ex}

\begin{ex}
	$\mathcal{O}=\{\emptyset,E\}$ est la \textbf{topologie grossière}. Nous verrons rapidement qu'elle est non métrisable dès que $E$ possède plus de deux éléments.
\end{ex}

\begin{dfn}
	Soit $(E,\mathcal{O})$ un espace topologique. On appelle fermé de $E$ toute partie $F$ de $E$ dont le complémentaire dans $E$ appartient à $\mathcal{O}$.
\end{dfn}

\begin{rmq}
	On vérifie facilement que l'ensemble des fermés est stable par intersection quelconque et par réunion finie.
\end{rmq}

\begin{dfn}
	Soit $a$ un élément de $E$, où $(E,\mathcal{O})$ est un espace topologique. On appelle voisinage de $A$ toute partie de $E$ telle qu'il existe $U\in\mathcal{O}$ vérifiant $a\in U$ et $U\subset A$.
\end{dfn}

\begin{rmq}
	On vérifie aisément que, comme pour les espaces métriques, que toute partie contenant un voisinage de $a$ est un voisinage de $a$, et qu'une intersection finie de voisinages de $a$ est un voisinage de $a$.
\end{rmq}

\begin{rmq}
	On définit de même les notions de points intérieurs et adhérents, ainsi que de densité, \textit{via} la notion de voisinage en adaptant les définitions équivalentes obtenues dans le cadre des espaces métriques.
\end{rmq}

\begin{dfn}
	Soit $(E,\mathcal{O})$ un espace topologique. On note $\mathcal{V}_a$ l'ensemble des voisinages de $a\in E$. On dit que la suite $(u_n)_{n\in\N}\in E^\N$ converge vers $a$ si : $$\forall V\in\mathcal{V}_a,\ \exists n_0\in\N\ : n\geqslant n_0\implies u_n\in V$$
\end{dfn}

\begin{rmq}
	On voit donc que la notion de convergence est proprement topologique et dépend fortement de la topologie considérée.
\end{rmq}

\begin{thm}{Métrique $\implies$ topologique}
	Soit $(E,d)$ un espace métrique. On note $\mathcal{O}$ l'ensemble des parties $U$ de $E$ telles que : $$\forall x\in U,\ \exists r>0,\ B_o(x,r)\subset U.$$ Alors $\mathcal{O}$ est une topologie, dite topologie associée à la distance $d$. Tout espace métrique est donc naturellement muni d'une structure d'espace topologique.
\end{thm}

\begin{rmq}
	Il s'agit donc bien d'une généralisation des sections précédentes.
\end{rmq}

\paragraph{Terminologie} Si une topologie $\mathcal{O}$ est la topologie associée à une distance $d$, on dit qu'elle est \textbf{métrisable.}

\begin{ex}
	La topologie grossière est non métrisable si $E$ possède plus de éléments. En effet, il est aisé qu'alors, toute suite d'éléments de $E$ converge vers tout élément de $E$. Si $E$ possède au moins deux éléments, il n'y a donc pas unicité de la limite, si bien que la topologie ne provient pas d'une distance.
\end{ex}

\begin{dfn}
	Soit $(E,\mathcal{O})$ un espace topologique et $A$ une partie de $E$. On note $$\mathcal{O}'=\{O\cap A\ \mid\ O\in\mathcal{O}\}.$$ Alors $\mathcal{O}'$ est une topologie sur $A$, appelée \textbf{topologie trace} de $\mathcal{O}$ sur $A$, qui munit $A$ d'une structure d'espace topologique.
\end{dfn}

\begin{rmq}
	Les notions d'ouverts, fermés et voisinages relatifs sont donc des notions topologiques liées à la topologie trace.
\end{rmq}

\section{Comparaison de normes}

Il est fréquent de rencontrer plusieurs normes sur un même espace vectoriel. une question apparaît alors : si une propriété (comme la convergence d'une suite ou le caractère borné, ouvert ou fermé d'une partie) est vraie pour une norme, l'est-elle également pour les autres ?\\
La réponse est négative dans le cas général, comme l'illustre l'exemple suivant.

\begin{ex}
	On considère la suite de fonctions $(f_n)$ définie sur $[0,1]$ par $f_n(x)=x^n$.
	\begin{enumerate}
		\item Un exemple précédent montrer que $(f_n)$ converge vers $0$ pour la norme un mais pas pour la norme infinie.
		\item En notant $g_n=(n+1)f_n$, on constate que la suite $(g_n)$ est bornée pour la norme un (tous ses éléments sont unitaires), mais pas pour la norme infinie.
	\end{enumerate}
\end{ex}

Cependant, nous allons voir ici que sous des condition supplémentaires, certaines propriétés sont conservées lors du passage d'une norme à une autre.

\subsection{Norme dominée par une autre}

\begin{dfn}
	Soit $N_1$ et $N_2$ deux normes sur $E$. On dit que $N_1$ est \textbf{dominée} par $N_2$, ou que $N_2$ \textbf{domine} $N_1$, s'il existe un réel $\alpha$ strictement positif tel que $N_1\leqslant\alpha N_2$, ou encore : $$\forall x\in E,\ N_1(x)\leqslant\alpha N_2(x).$$
\end{dfn}

\paragraph{Terminologie} Pour signifier que $N_1$ est dominée par $N_2$, on dit aussi que $N_2$ est \textbf{plus fine} que $N_1$.

\begin{rmq}
	La propriété $N_1\leqslant\alpha N_2$ s'interprète ainsi :
	\begin{itemize}
		\item pour tout $\varepsilon>0$, la boule pour la norme $N_2$ centrée en $0$ et de rayon $\varepsilon/\alpha$ est incluse dans la boule pour la norme $N_1$ de centre $0$ et de rayon $\varepsilon$.
		\item pour tout $\eta>à$, la boule pour la norme $N_2$ centrée en $0$ et de rayon $\eta$ est incluse dans la boule pour la norme $N_1$ de centre $0$ et de rayon $\eta\alpha$.
	\end{itemize}
\end{rmq}

Le résultat suivant assure que pour étudier le caractère dominé de la norme $N_1$ par la norme $N_2$, on peut se contenter de considérer les éléments $x\in E$ tels que $N_2(x)=1$.

\begin{thm}{}%prop
	Soit $N_1$ et $N_2$ deux normes sur $E$. La norme $N_1$ est dominée par la norme $N_2$ si, et seulement s'il existe un réel $\alpha$ strictement positif tel que : $$\forall x\in E,\ \lVert x\rVert=1\implies N_1(x)\leqslant\alpha.$$
\end{thm}
\begin{proof}
	Raisonnement par homogénéité.
\end{proof}

\begin{exo}
	La relation de domination pour les normes est-elle une relation d'ordre ?
\end{exo}

\begin{exo}
	Dans l'espace $\mathcal{C}([a,b],\K)$, montrer que les normes un et deux sont toutes les deux dominées par la norme infinie.
\end{exo}

\begin{thm}{}%prop
	Soit $N_1$ et $N_2$ deux normes sur $E$ telle que $N_1$ soit dominée par $N_2$. Si une partie de $E$ est bornée pour la norme $N_2$, alors elle l'est également pour la norme $N_1$.
\end{thm}

\begin{rmq}
	Avec les notations de la proposition précédente, si une application $f$ à valeurs dans $E$ (et donc en particulier une suite d'éléments de $E$) est bornée pour la norme $N_2$, alors elle l'est également pour la norme $N_1$.
\end{rmq}

\begin{thm}{}%prop
	Soit $N_1$ et $N_2$ deux normes sur $E$ telle que $N_1$ soit dominée par $N_2$. Si une suite converge vers un élément $l$ de $E$ pour la norme $N_2$, alors elle converge également vers $l$ pour la norme $N_1$.
\end{thm}

\begin{rmq}
	Étant données deux normes $N_1$ et $N_2$, pour montrer que $N_1$ n'est pas dominée par $N_2$, il suffit, d'après la proposition précédente, d'exhiber une suite qui converge pour la norme $N_2$ mais pas pour la norme $N_1$.
\end{rmq}

\begin{met}
	Dans la pratique, pour montrer que $N_1$ n'est pas dominée par $N_2$, on cherche souvent une suite qui, au choix : 
	\begin{itemize}
		\item tend vers $0$ pour $N_2$ mais pas pour $N_1$ ;
		\item est bornée pour $N_2$ mais qui ne l'est pas pour $N_1$.
	\end{itemize}
\end{met}

\subsection{Normes équivalentes}

\begin{dfn}
	Soit $N_1$ et $N_2$ deux normes sur $E$. On dit que $N_1$ et $N_2$ sont \textbf{équivalentes} s'il existe deux réels $\alpha$ et $\beta$ strictement positifs tels que : $$\forall x\in E,\ \alpha N_1(x)\leqslant N_2(x)\leqslant\beta N_1(x).$$
\end{dfn}

\begin{rmq}
	Deux normes sont donc équivalentes si, et seulement si, chacune est dominée par l'autre.
\end{rmq}

\begin{exo}
	Montrer que la relation d'équivalence des normes est une relation d'équivalence.
\end{exo}

\begin{exo}
	Montrer que pour $x\in\K^n$, on a : $$\lVert x\rVert_\infty\leqslant\lVert x\rVert_2\leqslant\lVert x\rVert_1\leqslant n\lVert x\rVert_\infty.$$ En déduire que ces trois normes sur $\K^n$ sont équivalentes.
\end{exo}

Beaucoup de propriétés sont conservées lorsqu'on passe à une autre qui lui est équivalente. \\
Donnons trois résultats, dont les deux premiers sont des conséquences immédiates de la remarque sur le fait que l'équivalence des normes correspond à une domination mutuelle et de résultats sur les normes qui en dominent d'autres.

\begin{thm}{Conservation du caractère borné d'une partie}%prop
	Soit $N_1$ et $N_2$ deux normes équivalentes sur $E$. une partie est bornée pour la norme $N_1$ si, et seulement si, elle l'est pour la norme $N_2$.
\end{thm}

\begin{rmq}
	Ce résultat s'applique en particulier aux applications bornées et aux suites bornées.
\end{rmq}

\begin{thm}{Conservation du caractère convergent d'une suite}
	Soit $N_1$ et $N_2$ deux normes équivalentes sur $E$. Une suite converge vers un élément $l$ de $E$ pour la norme $N_1$ si, et seulement si, elle converge vers $l$ pour la norme $N_2$.
\end{thm}

\begin{thm}{Conservation des ouverts et des fermés}%prop
	Soit $N_1$ et $N_2$ deux normes équivalentes sur $E$, et $A$ une partie de $E$. Alors :
	\begin{itemize}
		\item la partie $A$ est ouvert pour $N_1$ si, et seulement si, elle l'est pour $N_2$ ;
		\item la partie $A$ est fermée pour $N_1$ si, et seulement si, elle l'est pour $N_2$.
	\end{itemize}
\end{thm}

\begin{rmq}
	Les résultats précédents nous assurent que, pour étudier la convergence d'une suite ou encore le caractère borné, ouvert ou fermé (et donc, par extension, tout résultat reposant sur ces notions), si l'on dispose de plusieurs normes équivalentes, alors on peut choisir celle qui est la plus adaptée au problème.
\end{rmq}

\begin{met}
	En particulier, dans $\K^n$, on se demandera régulièrement quelle norme parmi la norme 1, la norme 2 et la norme infinie, est la plus adaptée à la situation.
\end{met}

\subsubsection*{Comparaison de normes}

\begin{met}
	\textbf{Comparer} des normes $N_1$ et $N_2$ signifie déterminer si $N_1$ (respectivement $N_2$) est dominée par $N_2$ (respectivement $N_1$).
\end{met}

\begin{exo}
	Soit $a$, $b$ et  $c$ trois réels vérifiant $a<b<c$. On considère les deux normes suivantes sur l'espace $\R[X]$ : $$N_1(P)=\int_{a}^{b}\lvert P\rvert\ \ \text{et}\ \ N_2(P)=\int_{a}^{c}\lvert P\rvert$$
	Comparer les deux normes $N_1$ et $N_2$.
\end{exo}

\subsubsection*{Équivalence des normes en dimension finie}

Le théorème fondamental suivant sera démontré ultérieurement.

\begin{thm}{}%thm
	Dans un espace de dimension finie, toutes les normes sont équivalentes.
\end{thm}


\end{document}